\chapter*{五年级数学知识点填空}
\section*{计算}
\subsection*{小数乘除法（五上）}
\subsubsection*{小数乘法}
\begin{enumerate}
    \item 小数乘法计算方法：先按\blkda{整数乘法}计算，再看\blkda[4]{因数中一共有几位小数}，就在积中从右往左数出几位点上小数点. 例如：$1.2 \times 0.04 = $\blkda{$0.048$}

    \item 补0规则：数位不够时，在整数积的\blkda{左}边补0，再\blkda{点小数点}。如 $0.03\times0.002$，整数积是6，总共有\blkda{$5$}位小数，需要补\blkda{$4$}个0，得到\blkda{$0.00006$}

    \item 估算检查：计算前先估算，防止小数点位置点错。例如，$0.3 \times 0.5$结果应该比$0.3$\blkda{小}，如果结果是$1.5$就明显错了.
\end{enumerate}

\subsubsection*{小数除法}
\begin{enumerate}
    \item 小数除法计算方法：
    \begin{dingautolist}{172}
        \item 除数是整数：先按\blkda[3]{整数除法}计算，再将商的小数点对齐，即\blkda[3]{商的小数点}要和\blkda[3]{被除数的小数点}对齐；如果有余数，在后面补0继续除，直到除尽或达到所需精度.
        \item 除数是小数：利用\blkda[3]{商不变性质}将除数变成整数.
    \end{dingautolist}
    例如：$4.5 \div 0.15 = $\blkda{30}

    \item 补0规则：当被除数移动小数点后位数不够时，在被除数\blkda{右}边用0补足。例如 $2.5 \div 0.005$，除数移\blkda{$3$}位变成\blkda{5}，被除数$2.5$移\blkda{$3$}位变成\blkda{$2500$}.
    
    \item 商的性质：
    \begin{dingautolist}{172}
        \item 一个正数（0除外）除以一个\blkda{小于}1的正数，商大于被除数。（如 $5 \div 0.5 = 10$）
        \item 一个正数除以一个\blkda{大于}1的数，商小于被除数。（如 $5 \div 2.5 = 2$）
    \end{dingautolist}
\end{enumerate}

\subsubsection*{循环小数}
\begin{enumerate}
    \item 循环小数是无限小数的一种（小数位数无限）。
    \item 循环小数是指在小数部分，从某一位起，\blkda[6]{一个或几个数字依次不断重复出现}的小数
    \item 重复出现的数字段称为\blkda{“循环节”}
    \item 表示法：为了简洁，我们只在\blkda[4]{循环节的首位和末位数字}上点一个点（循环点），或者用横线标出整个循环节
\end{enumerate}

\subsection*{负数（五下+拓展）}
\subsubsection*{负数的概念}
\begin{enumerate}
    \item 定义：比\blkda{0}小的数
    \item 表示法：用负号“$-$”表示，如：$-1$, $-2.5$, $-\dfrac{3}{4}$
    \item 意义/用途：\blkda[4]{表示与正数相反意义的量}，例如：
    \begin{dingautolist}{172}
        \item 温度：零下$5$度记作$-5$度
        \item 海拔：海平面以下$100$米记作 $-100$米
        \item 财务：欠债$100$元记作 $-100$元
    \end{dingautolist}
\end{enumerate}
\subsubsection*{数轴}
\begin{enumerate}
    \item 数轴三要素：
    \begin{dingautolist}{172}
        \item \blkda{原点}：参考起点
        \item \blkda{正方向}：通常（默认）向\blkda{右}
        \item \blkda{单位长度}：均匀的刻度间隔
    \end{dingautolist}
    \item 数轴上的比较规则：\blkda{右边}的数总是大于\blkda{左边}的数
    \item 相反数：一个数与其相反数关于\blkda{原点}对称
    \item 绝对值：\blkda[6]{一个数在数轴上对应点到原点的距离}，一个数和其相反数的绝对值\blkda{相等}，一个数的绝对值一定\blkda[2]{大于等于}0，即一个正数的绝对值等于\blkda{其本身}，一个负数的绝对值等于\blkda{其相反数}. 或者可以说，一个数有负号的时候，去掉负号就是其绝对值，没有负号的时候，其本身就是自己的绝对值.
\end{enumerate}
\subsubsection*{数的分类}
\begin{enumerate}
    \item 根据和0的大小关系可以将数分成：\blkda{正数}、\blkda{负数}和\blkda{0}.
    \item 根据是否有小数部分，可以将数分成：\blkda{整数}、\blkda{小数}.
    \item 从0开始递增的整数称作\blkda{自然数}.
    \item 可以表示成两个整数的分数的数称作是\blkda{有理数}，不能表示为两个整数的分数的数称作是\blkda{无理数}.
\end{enumerate}
\subsubsection*{负数运算规则}
\begin{enumerate}
    \item 大小比较：“负得越多”，值越小；绝对值越大的负数越\blkda{小}
    \item 加减法
    \begin{dingautolist}{172}
        \item 同号相加：绝对值\blkda[3]{相加}，符号\blkda{不变}，例如：$(+3) + (+5) = $\blkda{8}，$(-3) + (-5) = $\blkda{-8}
        \item 异号相加：绝对值\blkda[3]{相减（大减小）}，符号取绝对值\blkda{大}的数的符号，例如：$(+7) + (-3) = $\blkda{4}，$(-7) + (+3) = $\blkda{$-4$}
        \item 减去一个数等于加上\blkda[3]{它的相反数}，例如：$-5 - 3 = -5 + ($ \blkda{$-3$} $) =$ \blkda{$-8$}
    \end{dingautolist}
    \item 乘除法
    \begin{dingautolist}{172}
        \item 符号：同号得\blkda{正}，异号得\blkda{负}
        \item 绝对值：绝对值做相应的乘除法运算.
    \end{dingautolist}
    例如：
    \begin{dingautolist}{172}
        \item $(-3)\times 4 = $ = \blkda{$-12$}
        \item $(-3) \times (-4) = $ \blkda{$12$}
        \item $12 \div (-3) = $ \blkda{$-4$}
        \item $-12 \div (-3) = $ \blkda{$4$}
    \end{dingautolist}
\end{enumerate}
\subsection*{常见巧算方式}
\begin{enumerate}
    \item $2, 5$凑十
    
    $96 \times 0.125 = $ \blkda{$48$} $\times 0.25 = $ \blkda{$24$} $\times 0.5 = $ \blkda{$12$}
    \item 直接移动小数点
    \begin{dingautolist}{172}
        \item 乘10、100、1000…：小数点向\blkda{右}移动\blkda{1}位、\blkda{2}位、\blkda{3}位…
        \item 除以10、100、1000…：小数点向\blkda{左}移动\blkda{1}位、\blkda{2}位、\blkda{3}位…
    \end{dingautolist}

    \item 移动小数点提取公因式
    $1.3 \times 3 + 13 \times 1.7 = $ \blkda{$13$} $\times ($\blkda{$0.3$} $+1.7) = $\blkda{$26$} 
\end{enumerate}
\subsection*{运算律常见错误}
\begin{enumerate}
    \item $(2+3) \times 2 = $\blkda{2} $\times 2 + $\blkda{3} $\times 2$
    \item $(2-3) \times 2 = $\blkda{2} $\times 2$ \blkda{$-$} $3 \times 2$
    \item $(10+5) \div 5$ \blkda{$=$}$10 \div 5 + 5 \div 5$
    \item $15 \div (2+3)$ \blkda{$\neq$}$15 \div 2 + 15 \div 3$
\end{enumerate}
    
\section*{方程}
\subsection*{方程的概念与解方程（五上+五下）}
\begin{enumerate}
    \item 方程是含有\blkda{未知数}的\blkda{等式}
    \item 解方程的过程是对等式两边做同样的操作，将方程化简，得到未知数的值的过程

    ${3x} + 7 - 2\left( {4 - x}\right)  = {5x} - \left( {x + 9}\right)  + 3$

    去括号：\blkda[6]{${3x} + 7 - 8 + {2x} = {5x} - x - 9 + 3$}

    \blkda[4]{合并同类项}：${5x} - 1 = {4x} - 6$

    两边同时减去$4x$： \blkda[4]{$x-1 = -6$}

    \blkda[4]{两边同时$+1$}：$x = -5$
    
\end{enumerate}

\subsection*{列方程解应用题（五上+五下）}
解题步骤：
\begin{dingautolist}{172}
    \item \blkda[6]{找到未知量}（必然包含题目中问的量，可能存在多个未知量）
    \item \blkda[6]{寻找等量关系，列出方程}
    \item \blkda[6]{求解方程}
    \item \blkda[6]{回答问题}
\end{dingautolist}
两种不同的出题模式：
\begin{dingautolist}{172}
    \item 单一情景：给出某一个量的相关信息，可以通过计算得到该量，直接建立等式

    例如：哥哥弟弟年龄和的两倍等于爸爸的年龄，爸爸36岁，哥哥10岁，求弟弟年龄.

    \item 两个情景：给出两个情景，在每个情景中都可以计算得到同一个量，根据两个不同的情景建立等式

    例如：甲每小时生成若干个零件，计划6小时完成，甲提高效率后，每小时多生产2个零件，成提前1小时完成，求甲实际每小时生成多少个零件.
\end{dingautolist}
\subsubsection*{年龄问题}
情景核心：\blkda[12]{年龄差不变}
\subsubsection*{鸡兔同笼问题}
情景核心：\blkda[12]{设出“鸡”的数量，根据数量和表示出“兔”的数量}
\subsubsection*{工程问题}
情景核心：\blkda[12]{看好题干，厘清两个情景之间的关系}
\subsubsection*{行程问题}
情景核心：\blkda[12]{速度$\times$时间 $=$路程，区分清楚数字的具体意义，根据路程建立等式}
\section*{平均数}
\subsection*{平均数（五上）}
根据平均数$\times$\blkda{个数}$ = $\blkda{总数}，可以实现\blkda{总数}和\blkda{平均数}的转化
\section*{平面图形}
\subsection*{定义和性质}
\begin{enumerate}
    \item 平行四边形是\blkda[6]{两组对边分别平行的四边形}，两组对边\blkda{平行且相等}，对角相等，邻角互补，对角线互相平分
    \item 梯形是\blkda[6]{只有一组对边平行的四边形}，平行的两边叫做梯形的 \blkda{底}（上底、下底），不平行的两边叫做\blkda{腰}.有一个角（从而有两个角）是直角的梯形是\blkda{直角梯形}，两腰相等的梯形称作\blkda{等腰梯形}
\end{enumerate}

\subsection*{面积公式（五上）}
\begin{dingautolist}{172}
    \item 三角形面积公式：\blkda[4]{底$\times$高$\div 2$}
    \item 平行四边形面积公式：\blkda[4]{底$\times$高}
    \item 梯形形面积公式：\blkda[4]{（上底+下底）$\times$高$\div 2$}
\end{dingautolist}
\subsection*{割补技巧（五上）}
将不规则的，计算不了的图形的面积转化为规则的会计算面积的多边形

\section*{立体图形}
\subsection*{体积转换（五下）}
\begin{enumerate}
    \item $1 m^3 = $\blkda{1000}$dm^3 = $\blkda{$10^6$}$cm^3$
    \item $1L = 1$\blkda{$dm^3$}，$1mL = 1$\blkda{$cm^3$}
\end{enumerate}
\subsection*{长方体的基本性质（五下）}
\begin{enumerate}
    \item 长方体有\blkda{8}个顶点，\blkda{12}条棱，\blkda{6}个面，每个面都是\blkda{长方形}
    \item 立方体，即\blkda{正方体}，是所有棱长都\blkda{相等}的长方体
    \item 长方体的长是指\blkda[3]{向前的面的底边}，宽是指\blkda[3]{侧面的底边}，高是指\blkda{侧棱长}
    \item 长、宽、高分别为$a, b, c$的长方体的体积公式为\blkda{$V = abc$}
    \item 表面积是指\blkda[6]{立体图形所有的表面的面积之和}，长、宽、高分别为$a, b, c$的长方体的表面积公式为\blkda[4]{$S = 2(ab+bc+ca)$}
\end{enumerate}

\subsection*{长方体应用题常见条件}
\begin{enumerate}
    \item 两个一样的长方体拼到一起时，相当于减少了\blkda{2}个面的面积，若要得到表面积最大的情况，应选择\blkda{面积最小}的两个面拼合，若要得到表面积最小的情况，则相反.
    \item 当长方体的高变化时，\blkda{4}个面的面积发生了变化.
    \item 金属铸造，改变了物体的形状，但\blkda{体积}不发生变化.
    \item 常见的特殊的长方体模型包括：\blkda[6]{无盖鱼缸和通风管道}
\end{enumerate}

\section*{计数原理与概率}
\subsection*{计数原理（五下+拓展）}
\begin{enumerate}
    \item 乘法原理：完成一件事需要\blkda[3]{分多个步骤}，每个步骤有多种方法，且步骤之间相互独立。总方法数 = \blkda[4]{各步骤方法数的乘积}
    \item 加法原理：完成一件事有\blkda[3]{多类不同的方案}，每类方案都能独立完成任务，不同类方案之间不重叠。总方法数 = \blkda[4]{各类方法数的和}
    \item 排列：从$n$个不同元素中选出$m$个元素（$m \leq n$），并按照一定的顺序排成一列。强调\blkda[4]{顺序}（不同的顺序算不同的结果）。
    
    排列数$A_n^m =$\blkda[6]{$ n\times(n-1)\times\cdots \times(n-m+1)$}
    \item 组合：从$n$个不同元素中选出$m$个元素（$m \leq n$），不要求顺序。强调\blkda[4]{选出来即可，不排序}。
    
    组合数$C_n^m =$\blkda[6]{$\dfrac{A_n^m}{A_m^m} = \dfrac{n\times(n-1)\times\cdots \times(n-m+1)}{m\times(m-1)\times\cdots\times2\times1}$}，特别地，因为选出$m$个元素等价于不选$n-m$个元素，两者方案数相同，故可知：对于组合数满足$C_n^m=$\blkda{$C_n^{n-m}$}.
    \item 例子

    某班级有男生 4 人（记作 $M_1$, $M_2$, $M_3$, $M_4$），女生 3 人（记作 $F_1$, $F_2$, $F_3$）。
    \begin{dingautolist}{172}
        \item 从班级中选 1 名学生代表

        方案 A：从男生中选 1 人，有 4 种选法；方案 B：从女生中选 1 人，有 3 种选法。

        这是 分类（“或”关系），选男生 或 选女生即可完成任务，且类之间不重叠。用加法原理：$4+3=7$ 种。

        \item 选 1 名男班长和 1 名女班长（必须男女各一）
        
        分两步：1. 选男班长：4 种选法。2.选女班长：3 种选法。这两步都要完成（“且”关系），用乘法原理：$4 \times 3 = 12$ 种

        \item 从全班 7 人中选出 3 人，分别担任正班长、副班长、学习委员

        选出 3 人，还要分配不同职位，即要考虑顺序。$A_7^3 = 7 \times 6 \times 5 = 210$种.

        \item 从全班 7 人中选出 3 人作为代表去参加一个会议（代表之间无顺序差别）

        只选人不排序，是组合：$C_7^3 = \dfrac{7 \times 6 \times 5}{3 \times 2 \times 1} = 35$种.
    \end{dingautolist}
\end{enumerate}

\subsection*{古典概型（五下+拓展）}
\subsubsection*{概率}
概率就是一件事发生的\blkda{可能性}有多大，用一个从 $0$ 到 $1$ 之间的数表示。

\begin{enumerate}
    \item 如果概率是 $0$，表示\blkda{不可能发生}（比如：太阳从西边出来）。
    \item 如果概率是 $1$，表示\blkda{一定会发生}（比如：太阳每天都会升起）。
    \item 如果概率是 \blkda{$0.5$}，表示一半的可能性会发生（比如：抛硬币正面朝上）。
\end{enumerate}

\subsubsection*{事件}
\begin{enumerate}
    \item 随机事件
    
    “随机”就是结果不能\blkda{提前确定}，但所有可能的结果我们都知道。随机事件就是随机试验中可能发生也可能不发生的一件事情。
    
    \noindent 例子：
    \begin{enumerate}
        \item 抛一枚硬币，“正面朝上”是一个随机事件。
        \item 掷一个骰子，“点数是 4”是一个随机事件。
    \end{enumerate}

    \item 基本事件

    基本事件就是依某个角度观察（在某一个标准下）\blkda{不能再分}的最小结果.

    \noindent 例子（掷一个公平骰子）：
    \begin{enumerate}
        \item 当观察向上的点数的时候，基本事件有：点数是 1、点数是 2、点数是 3、点数是 4、点数是 5、点数是 6。
        \item 当观察向上的点数的奇偶性的时候，基本事件有： 奇数、 偶数。
    \end{enumerate}

    \item 事件
    
    “事件”可以很简单，只是一个\blkda{基本事件}，也可以由\blkda[3]{几个基本事件}组成。在概率里，事件就是我们关心的某种结果或结果组合。

    \noindent 例子（掷一个公平骰子）：
    \begin{enumerate}
        \item “点数是偶数”是一个事件，它包含 $ \{2,4,6\} $ 三种情况。
        \item “点数大于 4”是一个事件，它包含 $ \{5,6\} $ 两种情况。
    \end{enumerate}

    \item 样本空间就是\blkda[4]{所有基本事件的集合}，我们可以把它想象成一个“结果大全”

    \noindent 例子：
    \begin{enumerate}
        \item 抛硬币的样本空间：\{正面，反面\}。
        \item 掷骰子的样本空间：\{1, 2, 3, 4, 5, 6\}。
    \end{enumerate}
\end{enumerate}


\subsubsection*{古典概型}
\begin{enumerate}
    \item 古典概型是一种简单又公平的概率模型，它满足：
    \begin{dingautolist}{172}
        \item 试验只有\blkda{有限个}基本事件。
        \item 每个基本事件发生的可能性\blkda{完全相同}。
    \end{dingautolist}
    \noindent 例子：
    掷一个公平的骰子、抛一枚均匀的硬币、从一副洗匀的扑克牌中抽一张牌等。

    \item 古典概型的概率公式

    如果事件 $A$ 由一些基本事件组成，那么它的概率是：
    \blkda[6]{$P(A) = \dfrac{\text{事件 $A$ 包含的基本事件个数}}{\text{样本空间中基本事件的总数}}$}
    
    例子：掷骰子得偶数点的概率
    
    样本空间有 $6$ 个基本事件。\\
    事件 $A$ = “点数是偶数”包含 $ \{2,4,6\} $，共 $3$ 个基本事件。
    $P(A) = \dfrac{3}{6} = \dfrac{1}{2}$
\end{enumerate}